\documentclass[12pt]{article}
\usepackage[utf8]{inputenc}
\usepackage[T1]{fontenc}
\usepackage[margin=1in]{geometry}
\usepackage{setspace}
\usepackage{xcolor}
\usepackage[colorlinks=true, urlcolor=blue!60!black, linkcolor=blue!60!black]{hyperref}
\setlength{\parindent}{2em}
\setlength{\parskip}{1.25em}
\renewcommand{\baselinestretch}{1.0}
\title{\textit{New York Times Co. v. Sullivan}, 376 U.S. 254 (1964)}
\author{Case Brief}
\date{}
\begin{document}
\maketitle

% Case Brief: New York Times Co. v. Sullivan, 376 U.S. 254 (1964)
% Chapter: Free Speech & Press | Constitutional Law Briefing Book

\section{Case Metadata}
\begin{itemize}
  \item \textbf{Case Name:} \textit{\href{https://supreme.justia.com/cases/federal/us/376/254/}{New York Times Co. v. Sullivan}} (consolidated with \textit{Abernathy v. Sullivan}, No.~40)
  \item \textbf{Citation:} 376 U.S. 254, 84 S. Ct. 710, 11 L. Ed. 2d 686 (1964) (\href{https://www.law.cornell.edu/supremecourt/text/376/254}{Cornell LII}; \href{https://supreme.justia.com/cases/federal/us/376/254/}{Justia})
  \item \textbf{Court \& Dates:} Supreme Court of the United States; Argued January 6--7, 1964; Decided March 9, 1964
  \item \textbf{Vote:} 9--0 (unanimous judgment reversing and remanding the Supreme Court of Alabama, 273 Ala. 656, 144 So. 2d 25)
  \item \textbf{Majority Author:} Justice William J. Brennan, Jr. (joined by Chief Justice Warren and Justices Clark, Harlan, Stewart, and White)
  \item \textbf{Concurrences:} Justice Hugo L. Black (joined by Justice Douglas); Justice Arthur J. Goldberg, concurring in the result (joined by Justice Douglas)
  \item \textbf{Dissents:} None
\end{itemize}

\section{Facts and Procedural History}
On March 29, 1960, the \textit{New York Times} published a full-page editorial advertisement titled ``Heed Their Rising Voices,'' paid for at a cost of roughly \$4,800 by the Committee to Defend Martin Luther King and the Struggle for Freedom in the South (\href{https://supreme.justia.com/cases/federal/us/376/254/}{376 U.S. at 256--60}). The advertisement described a ``wave of terror'' against Southern civil rights demonstrators, appealed for funds to support student protests, voting rights, and Dr. Martin Luther King, Jr.'s defense against an Alabama perjury indictment, and appeared over the names of 64 prominent sponsors plus 20 Southern endorsers---including four African-American Alabama clergymen (among them Ralph D. Abernathy) who testified they never authorized use of their names (\href{https://www.law.cornell.edu/supremecourt/text/376/254}{376 U.S. at 257, 260}). Only about 394 copies of the issue circulated in Alabama, including roughly 35 in Montgomery County out of 650,000 nationwide (\href{https://supreme.justia.com/cases/federal/us/376/254/}{id. at 260 n.3}).

Respondent L.B. Sullivan, one of Montgomery's three elected City Commissioners and Commissioner of Public Affairs supervising the Police, Fire, Cemetery, and Scales Departments, was never named in the advertisement (\href{https://supreme.justia.com/cases/federal/us/376/254/}{id. at 256, 258}). He nevertheless sued the \textit{Times} and the four clergymen for libel based on paragraphs three and six, which contained conceded minor inaccuracies: student demonstrators on the Capitol steps sang the National Anthem rather than ``My Country, 'Tis of Thee''; nine students were expelled for a courthouse lunch-counter protest rather than the Capitol demonstration; most students boycotted classes for one day rather than the entire student body refusing to reregister; the Alabama State College dining hall was never padlocked; police were deployed near the campus on three occasions rather than ``ringing'' it; and Dr. King had been arrested four times rather than seven (with three arrests and both bombings of King's home predating Sullivan's tenure) (\href{https://www.law.cornell.edu/supremecourt/text/376/254}{id. at 258--59}).

As required for punitive damages under Ala. Code, Tit.~7, \S~914, Sullivan demanded a retraction; the \textit{Times} replied that it was ``puzzled'' as to how the text reflected on him and invited clarification, though it later retracted upon Alabama Governor John Patterson's demand (\href{https://supreme.justia.com/cases/federal/us/376/254/}{id. at 261--62}). Sullivan sued without replying and offered no proof of pecuniary loss. The Montgomery County Circuit Court instructed the jury that the statements were ``libelous per se,'' meaning falsity, malice, and general damages were presumed, and refused to require differentiation between compensatory and punitive damages (\href{https://supreme.justia.com/cases/federal/us/376/254/}{id. at 262}). The jury awarded the full \$500,000 sought against all defendants. The Supreme Court of Alabama affirmed, holding that the First Amendment does not protect libel, that the Fourteenth Amendment does not reach private civil suits, and that criticism of police automatically attaches ``of and concerning'' the supervising commissioner (\href{https://supreme.justia.com/cases/federal/us/376/254/}{id. at 263--64}). The U.S. Supreme Court granted certiorari (371 U.S. 946).

\section{Issue(s) Presented}
\begin{enumerate}
  \item Do the First and Fourteenth Amendments constrain a State's power to award civil libel damages in a lawsuit brought by a public official against critics of his official conduct?
  \item If so, what constitutional fault standard governs liability, and was the trial evidence constitutionally sufficient to establish either fault or that the statements were ``of and concerning'' Commissioner Sullivan?
\end{enumerate}

\section{Holding}
Yes. Reversing and remanding 9--0, the Court held that the First and Fourteenth Amendments prohibit a public official from recovering damages for a defamatory falsehood relating to his official conduct unless he proves with ``convincing clarity'' that the statement was made with \textbf{``actual malice''}---that is, with knowledge that it was false or with reckless disregard of whether it was false or not (\href{https://supreme.justia.com/cases/federal/us/376/254/}{376 U.S. at 279--80, 285--86}). Independently reviewing the record, the Court further held that the evidence was constitutionally insufficient both to show actual malice and to show that the impersonal criticism of governmental operations was ``of and concerning'' Sullivan (\href{https://www.law.cornell.edu/supremecourt/text/376/254}{id. at 285--92}).

\section{Reasoning}
Justice Brennan dismantled Alabama's framework in three steps. First, he rejected two threshold defenses: state-court enforcement of common-law tort rules constitutes Fourteenth Amendment ``state action,'' and an ``editorial advertisement'' protesting public abuses does not forfeit First Amendment protection merely because the newspaper was paid to publish it (\href{https://supreme.justia.com/cases/federal/us/376/254/}{376 U.S. at 265--66}, distinguishing \textit{Valentine v. Chrestensen}, 316 U.S. 52 (1942)).

Second, Brennan refused to grant the state-law label ``libel'' any ``talismanic immunity'' from First Amendment scrutiny (\href{https://supreme.justia.com/cases/federal/us/376/254/}{id. at 269}). Framing free debate on public issues as ``uninhibited, robust, and wide-open,'' he reasoned that erroneous statements are inevitable and need ``breathing space'' to survive (\href{https://supreme.justia.com/cases/federal/us/376/254/}{id. at 270--72}). Drawing on the ``central meaning of the First Amendment'' forged in the Sedition Act of 1798 controversy---whose invalidity had ``carried the day in the court of history''---Brennan explained that civil libel damages without criminal safeguards (here 1,000 times Alabama's maximum criminal libel fine, with no double-jeopardy bar against cascading suits) chill criticism even more severely than criminal sedition (\href{https://www.law.cornell.edu/supremecourt/text/376/254}{id. at 273--78}). Even a defense of truth fails because placing the burden of proving factual precision on critics induces self-censorship, causing speakers to ``steer far wider of the unlawful zone'' (\href{https://supreme.justia.com/cases/federal/us/376/254/}{id. at 278--79}). Analogizing to official immunity in \textit{Barr v. Matteo}, 360 U.S. 564 (1959), Brennan concluded that the ``citizen-critic'' deserves a reciprocal conditional privilege defeasible only by actual malice (\href{https://supreme.justia.com/cases/federal/us/376/254/}{id. at 282--83}).

Third, rather than merely remanding for new instructions because the general verdict masked presumed malice, the Court conducted an ``independent examination of the whole record'' to prevent First Amendment erosion (\href{https://supreme.justia.com/cases/federal/us/376/254/}{id. at 284--85}). Against the clergymen, there was zero proof of knowledge or recklessness. Against the \textit{Times}, failing to cross-check the ad against its own news morgue after relying on A. Philip Randolph's certification amounted at most to negligence, not reckless disregard (\href{https://supreme.justia.com/cases/federal/us/376/254/}{id. at 286--88}). Independently, Alabama's ``of and concerning'' inference impermissibly transmuted impersonal criticism of government into personal libel of officials---reviving seditious libel through ``legal alchemy'' (\href{https://www.law.cornell.edu/supremecourt/text/376/254}{id. at 288--92}).

\section{Concurring and Dissenting Opinions}
\begin{itemize}
  \item \textbf{Justice Black (concurring, joined by Justice Douglas):} Argued that the First and Fourteenth Amendments confer an \textit{absolute, unconditional immunity} to criticize public officials' conduct. Warning that 11 pending Alabama suits sought \$5.6 million against the \textit{Times} and \$1.7 million against CBS, Black dismissed ``actual malice'' as an ``elusive, abstract concept'' offering only ``evanescent protection'' against hostile local juries (\href{https://supreme.justia.com/cases/federal/us/376/254/}{376 U.S. at 293--97}).
  \item \textbf{Justice Goldberg (concurring in the result, joined by Justice Douglas):} Likewise urged an absolute privilege for criticism of official conduct, noting that juries probing subjective motivation can easily conflate falsity with malice, while carefully distinguishing purely private defamation of officials or private citizens as still actionable (\href{https://www.law.cornell.edu/supremecourt/text/376/254}{id. at 297--305}).
  \item \textbf{Dissents:} None.
\end{itemize}

\section{Key Quotations}
\begin{quote}
  ``Thus, we consider this case against the background of a profound national commitment to the principle that debate on public issues should be uninhibited, robust, and wide-open, and that it may well include vehement, caustic, and sometimes unpleasantly sharp attacks on government and public officials.'' \\
  --- Justice Brennan, \href{https://supreme.justia.com/cases/federal/us/376/254/}{376 U.S. at 270 (Justia)}
\end{quote}

\begin{quote}
  ``The constitutional guarantees require, we think, a federal rule that prohibits a public official from recovering damages for a defamatory falsehood relating to his official conduct unless he proves that the statement was made with `actual malice'---that is, with knowledge that it was false or with reckless disregard of whether it was false or not.'' \\
  --- Justice Brennan, \href{https://www.law.cornell.edu/supremecourt/text/376/254}{376 U.S. at 279--80 (Cornell LII)}
\end{quote}

\begin{quote}
  ``Whether or not a newspaper can survive a succession of such judgments, the pall of fear and timidity imposed upon those who would give voice to public criticism is an atmosphere in which the First Amendment freedoms cannot survive.'' \\
  --- Justice Brennan, \href{https://supreme.justia.com/cases/federal/us/376/254/}{376 U.S. at 278 (Justia)}
\end{quote}

\begin{quote}
  ``\,`Malice,' even as defined by the Court, is an elusive, abstract concept, hard to prove and hard to disprove. The requirement that malice be proved provides, at best, an evanescent protection for the right critically to discuss public affairs, and certainly does not measure up to the sturdy safeguard embodied in the First Amendment.'' \\
  --- Justice Black (concurring), \href{https://supreme.justia.com/cases/federal/us/376/254/}{376 U.S. at 293 (Justia)}
\end{quote}

\section{Doctrinal Significance and Modern Legacy (Through 2026)}
\textit{Sullivan} constitutionalized American defamation law and anchored modern First Amendment doctrine in democratic self-governance. The Court soon extended actual malice to ``public figures'' in \textit{Curtis Publishing Co. v. Butts}, 388 U.S. 130 (1967) (\href{https://www.law.cornell.edu/wex/defamation}{Cornell LII Wex}), drew the public/private boundary in \textit{\href{https://supreme.justia.com/cases/federal/us/418/323/}{Gertz v. Robert Welch, Inc.}}, 418 U.S. 323 (1974) (allowing states to set lower fault rules for private plaintiffs' actual damages while barring strict liability), applied \textit{Sullivan} to intentional infliction of emotional distress in \textit{\href{https://supreme.justia.com/cases/federal/us/485/46/}{Hustler Magazine, Inc. v. Falwell}}, 485 U.S. 46 (1988), and reinforced independent appellate review of actual malice in \textit{Bose Corp. v. Consumers Union}, 466 U.S. 485 (1984) (\href{https://www.law.cornell.edu/constitution-conan/amendment-1/defamation}{Constitution Annotated}) and \textit{\href{https://supreme.justia.com/cases/federal/us/491/657/}{Harte-Hanks Communications, Inc. v. Connaughton}}, 491 U.S. 657 (1989). In \textit{\href{https://www.supremecourt.gov/opinions/22pdf/600us1r51_g3bi.pdf}{Counterman v. Colorado}}, 600 U.S. 66 (2023), a 7--2 Court expressly drew on \textit{Sullivan}'s ``breathing space'' mens rea logic to require subjective recklessness in true-threats prosecutions.

Yet \textit{Sullivan} faces sustained originalist and institutional pressure through 2026. Justice Thomas has repeatedly urged overruling \textit{Sullivan} as a ``policy-driven'' departure from Founding-era common-law libel in \textit{\href{https://www.supremecourt.gov/opinions/18pdf/17-1542_ihdk.pdf}{McKee v. Cosby}}, 586 U.S. 1172 (2019), \textit{\href{https://www.supremecourt.gov/opinions/20pdf/20-1063_new_gfbi.pdf}{Berisha v. Lawson}}, 594 U.S. \_\_\_ (2021), \textit{\href{https://www.supremecourt.gov/opinions/23pdf/22-1125_c07d.pdf}{Blankenship v. NBCUniversal, LLC}}, 601 U.S. \_\_\_ (2023), and his \textit{Counterman} dissent. Justice Gorsuch joined the call in \textit{\href{https://www.supremecourt.gov/opinions/20pdf/20-1063_new_gfbi.pdf}{Berisha}}, questioning whether an actual-malice rule designed for 1964 print journalism subsidizes viral digital disinformation. Most recently, on June 29, 2026, Justice Thomas---joined by Justice Gorsuch---dissented from denial of certiorari in \textit{\href{https://www.supremecourt.gov/opinions/25pdf/25-770_2c8f.pdf}{Dershowitz v. Cable News Network, Inc.}}, 609 U.S. \_\_\_ (2026), again urging reconsideration. Still, the Court majority has consistently declined to grant review---including denying certiorari in \textit{Wynn v. Associated Press} in 2025 (\href{https://www.fire.org/news/blogs/ronald-kl-collins-first-amendment-news/trumps-executive-orders-due-process-breathtaking}{FIRE})---and \textit{Sullivan} continues to govern high-profile litigation, from the April 2025 retrial verdict for the \textit{Times} in \textit{\href{https://caselaw.findlaw.com/court/us-dis-crt-sd-new-yor/118036252.html}{Palin v. New York Times Co.}} to ongoing suits testing pleading thresholds and AI-era defamation (\href{https://www.lawcommentary.com/articles/trump-new-york-times-defamation-lawsuit-actual-malice}{Law Commentary}).

\section{Pedagogical Critiques and Exam Hypotheticals}
\textbf{Critiques:} Originalists (echoing Thomas) argue that neither the First nor Fourteenth Amendment displaced state civil libel law, noting Madison himself contemplated state remedies for injured reputations (\href{https://www.supremecourt.gov/opinions/18pdf/17-1542_ihdk.pdf}{McKee, slip op. at 12}). Conversely, absolutists (Black and Goldberg) contend that intrusive discovery into editorial states of mind chills underfunded speakers. Pragmatists add that in a low-cost digital ecosystem, actual malice may under-deter deliberate falsehoods while leaving wrongly maligned limited-purpose public figures without redress.

\textbf{Exam Hypothetical 1 (Standard Application):} A local investigative blog publishes a story accusing the City Police Chief of shredding body-camera footage, relying solely on one anonymous tipster whose prior tips proved fabricated, and deliberately declines to review publicly posted audit logs that would refute the claim. \textit{Analysis:} The Police Chief is a public official; while mere failure to investigate is negligence (\textit{Sullivan}), purposeful avoidance of the truth despite serious doubts about a known unreliable source can satisfy reckless disregard with convincing clarity (\textit{\href{https://supreme.justia.com/cases/federal/us/491/657/}{Harte-Hanks}}).

\textbf{Exam Hypothetical 2 (``Opposite Side'' Contrast):} Suppose instead that an unverified social-media post falsely accuses a private-citizen neighborhood mechanic---who never spoke to the press or entered any public debate---of overbilling city residents during a hurricane cleanup, and the local newspaper negligently repeats the claim without checking receipts. \textit{Contrast:} Unlike Commissioner Sullivan, the mechanic is a private figure who neither assumed the risk of public scrutiny nor commands media access to rebut falsehoods (\textit{\href{https://supreme.justia.com/cases/federal/us/418/323/}{Gertz}}, 418 U.S. at 344--45). Under \textit{Gertz}, the State may allow recovery of proven actual damages upon a showing of mere negligence, reserving \textit{Sullivan}'s actual-malice bar only for presumed or punitive damages on a matter of public concern.

\end{document}
